% Short paper on the Sphaera kernel -- working draft (structure and appendix only)
% Plan: about 40 pages; the Companion (concept DOI 10.5281/zenodo.19210728) stays the full record.
\documentclass[11pt]{article}
\usepackage[a4paper,margin=2.5cm]{geometry}
\usepackage[T1]{fontenc}
\usepackage{lmodern}
\usepackage{amsmath,amssymb,amsthm}
\usepackage{booktabs,longtable,array}
\usepackage{enumitem}
\usepackage[hidelinks]{hyperref}
\newtheorem{proposition}{Proposition}[section]
\newtheorem{lemma}[proposition]{Lemma}
\theoremstyle{definition}
\newtheorem{definition}[proposition]{Definition}
\theoremstyle{remark}
\newtheorem{remark}[proposition]{Remark}
\newcommand{\HH}{\mathbb H}
\newcommand{\OO}{\mathbb O}

\title{The Sphaera Kernel: a Biquaternion--Octonion Framework with Three Frames.\\
\large What holds, what is conditional, and what does not follow}
\author{Marc Brendecke}
\date{Working draft, \today}

\begin{document}
\maketitle

\begin{abstract}
% TODO (last): what is proved, what is conditional, what is withdrawn; the open problem of a common dynamics.
\end{abstract}

% Each section: content plan as comment; sources are labels in sources/research/ of the Companion.

\section{Introduction and the starting conjecture}\label{sec:intro}
% ~2 pp. Conjecture: matter, light and a tachyon sector are one object in three frames. Marked as conjecture.
% State up front: no derivation of physics from it is claimed; Sec. 8 lists what does not follow.

\section{The biquaternion kernel, the involution and the compensator}\label{sec:kernel}

This section collects what is proved about the basic object.  The results are elementary; they are stated with proofs
because the rest of the paper rests on them, and because an earlier version of the Companion read the kernel wrongly
(Remark~\ref{rem:signature}).  Layer numbers refer to the research rounds of the Companion~\cite{companion}.

\subsection{The kernel and its natural norm}

\begin{definition}[Biquaternion kernel]\label{def:kernel}
A \emph{kernel} is an element
\[
q = it\,\mathbf 1 + z_1e_1+z_2e_2+z_3e_3 \;\in\; \mathbb C\otimes_{\mathbb R}\HH,
\qquad t\in\mathbb R,\quad z=(z_1,z_2,z_3)\in\mathbb C^3,
\]
where \(e_1,e_2,e_3=e_1e_2\) are the imaginary quaternion units and \(i\) is the complex unit, which commutes with them.
The three coordinates \(z_k\) are called the \emph{kernels} \(K_1,K_2,K_3\).
\end{definition}

Let \(\bar q\) denote quaternion conjugation (\(e_k\mapsto-e_k\)) \emph{without} complex conjugation.

\begin{proposition}[Natural norm]\label{prop:norm}
The map \(N(q):=q\bar q\) takes values in \(\mathbb C\), is multiplicative, \(N(pq)=N(p)N(q)\), and
\[
N(q)=-t^2+s(z),\qquad s(z):=z_1^2+z_2^2+z_3^2 .
\]
\end{proposition}

\begin{proof}
The coefficients commute with the units, so the vector part of \(q\bar q\) cancels and \(e_k\bar e_k=1\) leaves
\((it)^2+\sum_kz_k^2\).  Multiplicativity is inherited from \(\HH\), since \(\overline{pq}=\bar q\,\bar p\) holds for the
complexified conjugation as well.
\end{proof}

\begin{remark}[The Hermitian norm has the wrong signature]\label{rem:signature}
The Companion originally used \(|q|^2=-t^2+\sum_k|z_k|^2\) and read it as a Minkowski metric.  As a real quadratic form in
\((t,\operatorname{Re}z,\operatorname{Im}z)\) it has signature \((1,6)\), and conjugation \(q\mapsto uqu^{-1}\) by real
quaternions acts by \(SO(3)\) on \(\operatorname{Re}z\) and \(\operatorname{Im}z\) separately, without boosts (layer
v4.8.0-r06).  The reading below replaces it.
\end{remark}

\subsection{The Lorentz group on the kernels}

\begin{proposition}[Reading B]\label{prop:readingB}
For \(u\in\mathbb C\otimes\HH\) with \(N(u)=1\), the map \(q\mapsto uqu^{-1}\) fixes \(t\) and acts on \(z\) by an element of
\(SO(3,\mathbb C)\); every element of \(SO(3,\mathbb C)\) arises.  Since \(SO(3,\mathbb C)\cong PSL(2,\mathbb C)\cong
SO^+(1,3)\), the kernels carry the proper orthochronous Lorentz group, with invariant \(s(z)\).
\end{proposition}

\begin{proof}
The unit group \(\{N(u)=1\}\) of \(\mathbb C\otimes\HH\cong M_2(\mathbb C)\) is \(SL(2,\mathbb C)\).  Conjugation fixes the scalar
part and preserves the complex three-dimensional space of pure biquaternions together with the restriction of \(N\), which is
\(s\).  This gives the double cover \(SL(2,\mathbb C)\to SO(3,\mathbb C)\).
\end{proof}

Writing \(z=E+iB\) with \(E,B\in\mathbb R^3\) identifies \(z\) with the Riemann--Silberstein vector of an electromagnetic
field~\cite{bialynicki}, and
\[
s=E^2-B^2+2i\,E\!\cdot\!B
\]
collects the two field invariants.  This is the \emph{working reading} of the kernel (``Reading~B''), adopted by the author
in layer v4.8.0-r06.  In it the scalar \(t\) is an invariant parameter, not a time coordinate; its role is open.

\begin{proposition}[Orbits]\label{prop:orbits}
The \(SO(3,\mathbb C)\)-orbits in \(\mathbb C^3\) are the level sets \(\{s=c\}\) for \(c\neq0\), the punctured null cone
\(\{s=0\}\setminus\{0\}\), and \(\{0\}\).
\end{proposition}

\begin{proof}
For \(s(z)=c\neq0\), the vector \(z/\sqrt c\) can be completed to a complex-orthonormal basis of determinant one, and the
rotation taking it to \(e_1\) maps \(z\) to \((\sqrt c,0,0)\).  For null vectors use Proposition~\ref{prop:nullcone}: the
spinor map is equivariant and \(SL(2,\mathbb C)\) acts transitively on \(\mathbb C^2\setminus\{0\}\).
\end{proof}

\subsection{The triple as a null configuration}

The starting point of the project was a triple of equal kernels at \(120^\circ\).  In Reading~B it has an invariant meaning.

\begin{lemma}[Triple lemma]\label{lem:triple}
For \(z\in\mathbb C^3\), with \(\omega=e^{2\pi i/3}\),
\[
z_1+z_2+z_3=0\ \text{ and }\ s(z)=0
\quad\Longleftrightarrow\quad
z=c\,(1,\omega,\omega^2)\ \text{up to order},\ c\in\mathbb C .
\]
\end{lemma}

\begin{proof}
With the elementary symmetric functions \(\varepsilon_1,\varepsilon_2,\varepsilon_3\) of the \(z_k\) one has
\(s=\varepsilon_1^2-2\varepsilon_2\).  So \(\varepsilon_1=s=0\) gives \(\varepsilon_2=0\), and the \(z_k\) are the roots of
\(\lambda^3-\varepsilon_3\), the three cube roots of \(\varepsilon_3\).  Conversely \(1+\omega+\omega^2=0=1+\omega^2+\omega^4\).
\end{proof}

The symmetric triple is therefore not only a phase condition: it is the Lorentz-invariant null condition together with
centring.  The null cone itself is classical:

\begin{proposition}[Null cone and spinors]\label{prop:nullcone}
For \(\psi\in\mathbb C^2\) put \(z(\psi)=(\psi_1^2-\psi_2^2,\;i(\psi_1^2+\psi_2^2),\;-2\psi_1\psi_2)\).  Then \(s(z(\psi))=0\),
every non-zero null vector has exactly the two preimages \(\pm\psi\), the map intertwines \(SL(2,\mathbb C)\) on \(\psi\) with
\(SO(3,\mathbb C)\) on \(z\), and the projectivized null cone is a smooth conic in \(\mathbb CP^2\), isomorphic to
\(\mathbb CP^1\cong S^2\).
\end{proposition}

\begin{proof}
With \(a=\psi_1^2\), \(b=\psi_2^2\): \(s=(a-b)^2-(a+b)^2+4ab=0\).  For null \(z\), \(\psi_1^2=\tfrac12(z_1-iz_2)\) and
\(\psi_2^2=-\tfrac12(z_1+iz_2)\), and \(z_3^2=4\psi_1^2\psi_2^2\) fixes the relative sign.  The conic is the Veronese image of
\(\mathbb CP^1\).
\end{proof}

\subsection{Sectors and the duality rotation}

\begin{definition}[Sectors; working assignment]\label{def:sectors}
A kernel with \(s(z)>0\) is called of \emph{matter type}, with \(s(z)<0\) of \emph{tachyon type}, and with \(s(z)=0\)
\emph{light-like}.
\end{definition}

This assignment is the author's working choice (layer v4.8.0-r07).  The results below are stated for the sign of \(s\) and do
not depend on it.  That the sectors \emph{are} matter, tachyon and light is the starting conjecture, not a result.

The \emph{duality rotation} \(z\mapsto e^{i\alpha}z\) rotates \(E\) into \(B\) and multiplies \(s\) by \(e^{2i\alpha}\).  The
quarter turn \(z\mapsto iz\) maps \(s\mapsto-s\); it exchanges the two massive types, fixes the null cone, and is the map that
takes a matter kernel to its tachyon image (layer v4.9.0-r28).

\begin{proposition}[Duality angle]\label{prop:duality}
Every \(z\) with \(s(z)\neq0\) can be written as \(z=e^{i\alpha}z_m\) with \(s(z_m)=|s(z)|>0\); here
\(\alpha=\tfrac12\arg s(z)\) is unique modulo \(\pi\) and \(z_m\) is unique up to sign.  Under
\(SO(3,\mathbb C)\times U(1)_{\rm dual}\), acting by \(z\mapsto e^{i\alpha}Rz\), the orbits in \(\mathbb C^3\setminus\{0\}\) are
the shells \(\{|s|=r\}\), \(r>0\), and the punctured null cone.
\end{proposition}

\begin{proof}
\(s(e^{i\alpha}z_m)=e^{2i\alpha}s(z_m)\), so \(s(z_m)\) is real positive exactly when \(2\alpha\equiv\arg s\pmod{2\pi}\).  The
orbit statement follows from Proposition~\ref{prop:orbits}, since the duality circle runs through all values of \(s\) of fixed
modulus and fixes \(s=0\).
\end{proof}

Matter-type and tachyon-type kernels of equal \(|s|\) thus lie in one orbit, related by a quarter turn, and the light-like
kernels form the only orbit fixed by the whole duality circle.  This is the precise form that the starting conjecture --- one
object in three frames --- takes on the kernel: two massive types exchanged by duality, and a duality-invariant fix-sector.
Whether the duality angle is physical or a gauge freedom is not decided.

\subsection{The involution and the compensator}

\begin{definition}[Involution and compensator]\label{def:compensator}
Let \(J_B z:=i\bar z\), with complex conjugation of the components.  Then \(J_B\) is antilinear and \(J_B^2=1\).  The
\emph{Heisenberg compensator} and its partner are
\[
C:=\tfrac12(1+J_B),\qquad A:=\tfrac12(1-J_B).
\]
\end{definition}

\begin{proposition}[Two ports]\label{prop:ports}
\(C\) and \(A\) are real-linear, \(C+A=1\), \(C^2=C\), \(A^2=A\), \(CA=AC=0\), and for the real inner product
\(\operatorname{Re}\langle Cf,Af\rangle=0\), so \(|Cf|^2+|Af|^2=|f|^2\).  Moreover \(A(f)=-i\,C(if)\): the two ranges differ by
a duality quarter turn.
\end{proposition}

\begin{proof}
The algebraic identities follow from \(J_B^2=1\).  Since \(|J_Bf|=|f|\),
\(\operatorname{Re}\langle Cf,Af\rangle=\tfrac14\operatorname{Re}\bigl(\langle J_Bf,f\rangle-\langle f,J_Bf\rangle\bigr)=0\), and the
cross terms \(\pm\tfrac12\operatorname{Re}\langle f,J_Bf\rangle\) cancel in \(|Cf|^2+|Af|^2=\tfrac12(|f|^2+|J_Bf|^2)\).
Finally \(C(if)=\tfrac12(if+J_B(if))=\tfrac12(if-iJ_Bf)=iA(f)\).
\end{proof}

The author's picture was a two-arm interferometer: one arm carries \(f\), the other \(J_Bf\), and the two output ports carry
\(Cf\) and \(Af\) (layer v4.9.0-r03).  Since \(J_B\) is antilinear, \(C\) has no complex matrix.  It becomes an ordinary
matrix on the real \emph{doublet} \(F=[\,E\mid B\,]\in\mathbb R^{3\times2}\):

\begin{proposition}[Doublet form]\label{prop:doublet}
On \(F\), acting from the right, \(J_B\) is \(\sigma_x\), the duality rotation by \(\alpha\) is the rotation \(R_\alpha^{\mathsf T}\),
and complex conjugation is \(\sigma_z\); they generate \(O(2)\).  The compensator is the orthogonal projection
\(C_{\rm HC}=\tfrac12(1+\sigma_x)\).  The Gram matrix \(G=F^{\mathsf T}F\) is
\[
G=u\,\mathbf 1+\tfrac12\bigl(\operatorname{Re}s\;\sigma_z+\operatorname{Im}s\;\sigma_x\bigr),\qquad
u=\tfrac12(E^2+B^2),
\]
and with \(b=(\operatorname{Re}s,\operatorname{Im}s)/(2u)\) the normalized state \(\rho=G/\operatorname{tr}G\) lies in the unit disk.
Matter-type fields have \(b_1>0\), tachyon-type \(b_1<0\), null fields \(b=0\).  \(J_B\) reflects \(b_1\), while the quarter
turn \(z\mapsto iz\) reverses the whole vector \(b\).
\end{proposition}

\begin{proof}
\(i(E-iB)=B+iE\) exchanges the columns; \(e^{i\alpha}(E+iB)\) rotates them.  \(G_{11}=E^2\), \(G_{22}=B^2\),
\(G_{12}=E\cdot B\).  A duality rotation conjugates \(G\) by \(R_\alpha\), which turns the traceless part by \(2\alpha\).
\end{proof}

Both \(J_B\) and the quarter turn send matter-type to tachyon-type kernels; they differ on the second invariant
\(\operatorname{Im}s=2E\cdot B\).  Three facts are worth stating plainly.  The compensator is a projection, not a dynamical
law.  Its two ports are real and imaginary parts of the same field, not two independent sectors (Proposition~\ref{prop:ports}).
And nothing in this section assigns a field equation to \(z\); it is kinematics on one point.

\medskip\noindent\emph{Status.}  Propositions~\ref{prop:norm}--\ref{prop:doublet} and Lemma~\ref{lem:triple} are proved.
Definition~\ref{def:sectors} is a working choice; its physical reading is the starting conjecture.  The role of \(t\) is open.

\section{Triples, the Hopf cascade, octonions and the compensator matrix}\label{sec:octonions}

The kernel of Section~\ref{sec:kernel} is the first of several layers.  This section states the structures that the Companion
builds on top of it, and separates what is classical mathematics from what is identification.

\subsection{Three meanings of ``triple''}

The word \emph{triple} (German \emph{Triole}) names three different structures in the Companion (layer v4.9.0-r37).  Several
errors found in earlier versions came from identifying them.
\begin{enumerate}[label=(\roman*)]
\item The \emph{frame triple}: matter, light and tachyon as three readings of one object, with light as the fix-sector of the
exchange of the other two (Proposition~\ref{prop:duality}).
\item The \emph{kernel triple}: three kernels of equal modulus at \(120^\circ\), \(z=c(1,\omega,\omega^2)\), which is the centred
null configuration of Lemma~\ref{lem:triple}.  Inside the octonions with one imaginary unit fixed, such a \(\mathbb Z_3\) is a colour
rotation (layer v4.9.0-r36).
\item The \emph{slot triple}: the three arguments of the triality form \(t(x,y,z)=\operatorname{Re}((xy)z)\) on \(\OO^3\).
\end{enumerate}

\begin{proposition}[Triality]\label{prop:triality}
Let \(\mathfrak g=\{(A,B,C)\in\mathfrak{so}(8)^3:\ t(Ax,y,z)+t(x,By,z)+t(x,y,Cz)=0\}\) and \(\tau(x,y,z)=(y,z,x)\).  Then
\(\dim\mathfrak g=28\), each projection \(\mathfrak g\to\mathfrak{so}(8)\) is an isomorphism, and the three slot representations are
the three inequivalent eight-dimensional representations of \(\mathrm{Spin}(8)\).  The shift \(\tau\) preserves \(t\) and induces an
outer automorphism of order three; its fixed elements \((A,A,A)\) form \(\operatorname{Der}(\OO)=\mathfrak g_2\), and those that
also annihilate \(e_7\) form \(\mathfrak{su}(3)\).
\end{proposition}

\begin{proof}
This is the classical triality of \(\mathrm{Spin}(8)\)~\cite{baez}.  The invariance \(t(x,y,z)=t(y,z,x)\) follows from
\(\operatorname{Re}((xy)z)=\operatorname{Re}(x(yz))\) and \(\operatorname{Re}(ab)=\operatorname{Re}(ba)\).  The dimensions were also checked by
solving the linear conditions numerically (layer v4.9.0-r37).
\end{proof}

In this picture a generation index is what the slot shift permutes and colour is what it fixes: orbit and fixed group of one
motion.  That the three slots \emph{are} the three physical generations is not shown.  Nor does the structure forbid a fourth
generation; it only excludes a fourth inequivalent eight-dimensional representation (\(\operatorname{Out}(\mathrm{Spin}(8))=S_3\)).

\subsection{The Hopf cascade}

The three classical Hopf fibrations~\cite{hopf1931,hopf1935}
\[
S^1\to S^3\to S^2=\mathbb CP^1,\qquad
S^3\to S^7\to S^4=\HH P^1,\qquad
S^7\to S^{15}\to S^8=\OO P^1
\]
organize the layers of the Companion.  The conjectural identification is: the base \(S^2\) is the light layer (Weyl spinors), \(S^4\)
the matter layer (Dirac spinors), and \(S^8\) the layer of the three generations (octonionic spinors), with the fibres
\(S^1\cong U(1)\) and \(S^3\cong SU(2)\) as gauge groups.

Two points of this reading are now fixed by proof, and two corrections are recorded.
\begin{itemize}
\item In Reading~B the directions of light-like kernels form exactly \(\mathbb CP^1\) (Proposition~\ref{prop:nullcone}).  So the
first base is not an analogy: it is the projectivized null cone of the kernel.
\item A normalized spinor \(\psi\) is a qubit; its Bloch vector is the Hopf image \(S^3\to S^2\), and
\(k(\psi)=(|\psi|^2,\psi^\dagger\sigma\psi)\) is a light-like four-vector.  Every future time-like vector is a sum
\(k(\psi_a)+k(\psi_b)\) (layer v4.8.0-r06).  This matches single spinors as light-like and pairs as massive readouts; it does not
derive space-time points.
\item The fibre \(S^7\) is not a group.  \(G_2\) is the automorphism group of \(\OO\), and \(SU(3)\) the stabilizer of an imaginary
unit; an earlier annotation of \(S^7\) as a group containing \(G_2\) was wrong (4.10 consolidation).
\item The chain \(S^{15}\to S^8\to S^4\to S^2\) is not a composition of Hopf maps (layer v4.8.0-r09).  The three fibrations stand
side by side; a cascade operator that runs through all three is a separate assumption.
\end{itemize}

\subsection{The compensator matrix}

The Companion asks for one operator that carries an object through all layers and contains the compensator.  It was called
the Millennium matrix; since layers v4.9.0-r75/r76 it is called the \emph{compensator matrix} \(\mathcal M_{16}\).  What the
Companion actually requires of it are the following identities on a common carrier.

\begin{definition}[Compensator-matrix conditions]\label{def:cm}
Operators \(\mathcal J,\mathcal C,\mathcal M\) on a real or complex vector space satisfy the \emph{compensator-matrix conditions} if
\[
\mathcal J^2=1,\qquad \mathcal C^2=\mathcal C,\qquad
\mathcal M\mathcal J=\mathcal J\mathcal M,\qquad \mathcal M\mathcal C=\mathcal C\mathcal M .
\]
\end{definition}

The Companion states that under these conditions \(\mathcal C\) is the spectral projector of \(\mathcal M\) onto its closure-stable
eigenspace.  Commutation alone does not give this: \(M=\operatorname{diag}(1,2,3)\) commutes with \(C=\operatorname{diag}(1,1,0)\)
without being scalar on \(\operatorname{im}C\), and \(M=\operatorname{diag}(2,2,3)\), \(C=\operatorname{diag}(1,0,0)\) shows that even
scalar action misses part of the eigenspace.  A correct statement needs two further hypotheses:

\begin{proposition}[Conditional spectral identification]\label{prop:spectral}
Let \(M\) be a bounded self-adjoint operator on a complex Hilbert space and \(C\neq0\) an orthogonal projector with
\[
MC=CM=\lambda_0C\quad\text{and}\quad\ker(M-\lambda_0)\cap\ker C=\{0\}
\]
for some real \(\lambda_0\).  Then \(\operatorname{im}C=\ker(M-\lambda_0)\), and \(C\) is the orthogonal eigenprojection of \(M\) for
\(\lambda_0\).
\end{proposition}

\begin{proof}
\(MC=\lambda_0C\) gives \(\operatorname{im}C\subset\ker(M-\lambda_0)\).  If \(Mx=\lambda_0x\), then \((1-C)x\) lies in
\(\ker(M-\lambda_0)\) by commutation and in \(\ker C\), hence vanishes, so \(x=Cx\).
\end{proof}

The hypotheses are assumptions, not consequences of the kernel.  The only place where the conditions have been solved
completely is the doublet of Proposition~\ref{prop:doublet}:

\begin{proposition}[Compensator matrix on the doublet]\label{prop:cm-doublet}
With \(\mathcal J=\sigma_x\) and \(\mathcal C=C_{\rm HC}\) acting on \(F=[\,E\mid B\,]\) from the right, the real \(2\times2\) matrices
satisfying Definition~\ref{def:cm} are exactly
\[
M=\lambda_+C_{\rm HC}+\lambda_-A,\qquad \lambda_\pm\in\mathbb R .
\]
For \(\lambda_+\neq\lambda_-\), \(C_{\rm HC}\) is the spectral projector of \(M\) for \(\lambda_+\).  The members of determinant one are
\(e^{\theta\sigma_x}=e^{\theta}C_{\rm HC}+e^{-\theta}A\); they preserve \(\operatorname{Re}s=E^2-B^2\).  Spatial rotations (acting from the
left) commute with all of them.  Boosts and duality rotations other than by \(0\) and \(\pi\) violate the conditions.
\end{proposition}

\begin{proof}
\(M\sigma_x=\sigma_xM\) is a linear condition whose solutions are spanned by \(1\) and \(\sigma_x\), equivalently by
\(C_{\rm HC}\) and \(A\).  \(R_\alpha\) commutes with \(\sigma_x\) only for \(\alpha\in\{0,\pi\}\), and a boost couples \(E\) and \(B\)
through cross products and is not of the form \(F\mapsto LFM\) (layers v4.9.0-r04, r05).
\end{proof}

On the doublet the compensator matrix can therefore only weight the two ports of the compensator; it is not the Lorentz transport
of the kernel.  This restriction is called the \emph{port-weighting matrix}.  The full \(\mathcal M_{16}\) of the Companion, with
its index \(16\) and its decomposition into four layer terms, has no realization so far.  One partial realization exists on
\(\mathbb C\otimes\HH\otimes\OO\): there a compensator of the same form, \(C_{\rm ch}=\tfrac12(1+J)\) with \(J=iL_{e_7}\), acts as the
chirality projector of one generation (Section~\ref{sec:generation}).  That one generation with a right-handed neutrino has \(16\)
Weyl states is recorded, not interpreted.

\medskip\noindent\emph{Status.}  Propositions~\ref{prop:triality}, \ref{prop:spectral} and \ref{prop:cm-doublet} are proved; the
first is classical.  The identification of the Hopf bases with light, matter and generations is a conjecture, proved only for the
first base in the sense of Proposition~\ref{prop:nullcone}.  The compensator matrix is defined by conditions; a realization beyond
the doublet and the chirality projector is open.


\section{One generation from \texorpdfstring{$\mathbb C\otimes\HH\otimes\OO$}{C x H x O}}\label{sec:generation}
% ~6 pp. r34, r35, r41-r43; conditional on named choices; comparison with Dixon, Furey, Gresnigt. QAM row 6.

\section{The frame offset \texorpdfstring{$\beta_{\rm sym}$}{beta\_sym}}\label{sec:beta}
% ~3 pp. thm:symmetric-frame; counting reading (r63) and its limit (r68, SU(2) vs U(1)). QAM row 7.

\section{Readouts and pre-registered predictions}\label{sec:readouts}
% ~5 pp. Omega_M (r67, present epoch only); Brannen/Koide form (r10, r46); H9, H15, H42, H46 with criteria. QAM row 8.

\section{The OP5 comparison model}\label{sec:op5}
% ~5 pp. r69-r74: theta gauge symmetry, Cartan test, four hidden sixth-order conditions, precision limit;
% parallel runs (p <= 12). Comparison model only; OP5 not closed. Candidate for a separate paper.

\section{What does not follow}\label{sec:not}
% ~3 pp. Withdrawn: Riemann-based steps (r04), GUE derivation (r05, r08), old Lambda chain, frame-shifted constants (r64),
% P1-P6; dependency map subsec:v480-r04-dependency-map. Status tags [M]/[T]/[P] and theorem/conditional/reading/open.

\section{Open problems: what a common dynamics would have to do}\label{sec:open}
% ~3 pp. OP1-OP5 short; the seven requirements from wiki "Offene Probleme"; H48.

\appendix
\section{Method}\label{app:method}
% ~1 p. Pre-registration, guards, Lean formalisation, independent recomputation.

% Appendix: the key statements in standard notation and in QAM notation (variant 1, author decision 2026-11-03)
\section{The key statements in QAM notation}
\label{app:qam}

The Companion~\cite{companion} records its objects in a coded language, QAM, which sits on ordinary Zermelo--Fraenkel set
theory.  It adds no axioms.  Eight tagged tuples recognise the kinds of object, and the first entry of each tuple is a type tag,
so no object belongs to two kinds:
\[
\begin{array}{llll}
\text{state} & \langle 0,a\rangle, &
\text{layer} & \langle 1,\ell,D\rangle,\\
\text{operator} & \langle 2,\ell_1,\ell_2,G\rangle, &
\text{readout} & \langle 3,\ell,W,R\rangle,\\
\text{branch} & \langle 4,b,B\rangle, &
\text{shell} & \langle 5,r,s,T\rangle,\\
\text{witness} & \langle 6,u\rangle, &
\text{tensor} & \langle 7,p,q\rangle .
\end{array}
\]
Here \(D\) is the carrier of a layer, \(G\) the graph of an operator from layer \(\ell_1\) to layer \(\ell_2\), \(W\) the set of
readout values and \(R\) the readout rule.  Three relations are used below: \(\mathrm{QAdmTransport}(o,s,t)\) (the operator
\(o\) carries the state \(s\) admissibly to \(t\)), \(\mathrm{QRead}(r,s,w)\) (the readout \(r\) reads \(s\) as the value
\(w\)) and \(\mathrm{QTensorAdm}(x,y)\) (the composite \(x\otimes y\) is admitted).  In the two-sided display
\({}^{A}{}_{B}\mathcal T^{i}{}_{j}\) the front indices carry the frame or readout side and the back indices the structure
side.  Only upper--lower pairs contract; two upper or two lower indices may be compared only through a declared comparison
tensor.  Full definitions are in the QAM appendix of the Companion.

The table gives each key statement of this paper in both notations.  The QAM column is not shorter.  It records what the
standard formula leaves to the surrounding text: on which layer an object lives, in which frame a quantity is read, whether a
value is derived or put in, and which choices were made.  The encodings were written for this paper from the QAM primitives;
the Companion does not contain them in this form.

\begin{center}
\small
\renewcommand{\arraystretch}{1.35}
\begin{longtable}{@{}>{\raggedright\arraybackslash}p{0.27\textwidth}>{\raggedright\arraybackslash}p{0.37\textwidth}>{\raggedright\arraybackslash}p{0.28\textwidth}@{}}
\toprule
standard notation & QAM notation & what the QAM form adds\\
\midrule
\endfirsthead
\toprule
standard notation & QAM notation & what the QAM form adds\\
\midrule
\endhead
\multicolumn{3}{@{}l}{\textbf{1. Biquaternion kernel} (\S\ref{sec:kernel})}\\
\(q=it\,\mathbf 1+z_1e_1+z_2e_2+z_3e_3\), \(z_k\in\mathbb C\), \(e_3=e_1e_2\)
& \(\ell_n=\langle 1,n,\HH\otimes\mathbb C\rangle\), \(s_q=\langle 0,q\rangle\) with \(q\in D(\ell_n)\)
& the stratum \(n\) on which the state lives\\[0.3em]
\multicolumn{3}{@{}l}{\textbf{2. Heisenberg compensator} (\S\ref{sec:kernel})}\\
\(C(f)=\tfrac12(f+Jf)\), \(J^2=1\), \(C^2=C\)
& \(o_C=\langle 2,\ell,\ell,G_C\rangle\), \(G_C=\{(\langle0,f\rangle,\langle0,\tfrac12(f+Jf)\rangle)\}\);
\(\mathrm{QAdmTransport}(o_C,s,o_C(s))\), \(G_C\circ G_C=G_C\)
& \(C\) is an operator within one layer, not a readout; idempotence is a statement about its graph\\[0.3em]
\multicolumn{3}{@{}l}{\textbf{3. Sector rule} (\S\ref{sec:kernel})}\\
\(\sigma(q)=z_1^2+z_2^2+z_3^2\) (\(s\) in the Companion): matter \(\sigma>0\), tachyon \(\sigma<0\), light \(\sigma=0\)
& \(r_\sigma=\langle 3,\ell,\{+,0,-\},R_\sigma\rangle\), \(R_\sigma(\langle0,q\rangle)=\operatorname{sgn}\sigma(q)\);
\(\mathrm{QRead}(r_\sigma,s_q,+)\;\Leftrightarrow\;\)matter
& the sector is a readout of the state, not a second kind of object\\[0.3em]
\multicolumn{3}{@{}l}{\textbf{4. Frame change matter \(\to\) tachyon} (\S\ref{sec:kernel})}\\
\(z\mapsto iz\), hence \(\sigma\mapsto-\sigma\)
& \({}^{T}{}_{M}\mathcal J^{i}{}_{j}\,z^{j}=(iz)^{i}\); a matter quantity \({}^{M}X\) and a tachyon quantity \({}^{T}Y\) are
compared only as \(G_{MT}\,{}^{M}X\,{}^{T}Y\) with a declared \(G_{MT}\)
& source and target frame in the indices; raw comparison across frames is not well-formed\\[0.3em]
\multicolumn{3}{@{}l}{\textbf{5. Compensator matrix} (\S\ref{sec:octonions})}\\
\([\mathcal M_{16},J]=[\mathcal M_{16},C]=0\); \(C\) is a spectral projector of \(\mathcal M_{16}\)
& \(o_{\mathcal M}=\langle 2,\ell,\ell,G_{\mathcal M}\rangle\) with \(G_{\mathcal M}\circ G_J=G_J\circ G_{\mathcal M}\) and
\(G_{\mathcal M}\circ G_C=G_C\circ G_{\mathcal M}\)
& nothing beyond the standard form; listed for completeness\\[0.3em]
\multicolumn{3}{@{}l}{\textbf{6. One generation from \(\mathbb C\otimes\HH\otimes\OO\)} (\S\ref{sec:generation})}\\
\(C_{\rm ch}=\tfrac12(1+J)\), \(J=iL_{e_7}\); \(Y/2=Q-\tfrac12+(1-C_{\rm ch})T_3\) gives the 16 left-handed Weyl states with \(\nu_R\)
& \(x=\langle 7,\mathbb C\otimes\HH,\OO\rangle\), \(\mathrm{QTensorAdm}(\mathbb C\otimes\HH,\OO)\); the choice
\(J=iL_{e_7}\) is a branch \(\langle 4,b_J,B\rangle\); charges are readouts \(\langle 3,\ell,\{0,\tfrac13,\tfrac23,1\},R_Q\rangle\)
& the choice of \(J\) is an explicit branch, so the construction is visibly conditional\\[0.3em]
\multicolumn{3}{@{}l}{\textbf{7. Frame offset} (\S\ref{sec:beta})}\\
\(\beta_{\rm sym}=\sqrt{(\gamma_L-1)/(\gamma_L+1)}\approx0.549\), \(\gamma_L=\gamma_{\rm BI}/\gamma_{\rm string}\)
& \({}^{M}\beta=R_\beta(u)\) with the witness \(u=\langle 6,(\gamma_{\rm BI},\gamma_{\rm string})\rangle\)
& the value rests on an input witness (two horizon countings), not on a derivation, and is read in the matter frame\\[0.3em]
\multicolumn{3}{@{}l}{\textbf{8. Matter fraction} (\S\ref{sec:readouts})}\\
\(\Omega_M=\beta_{\rm sym}^2\approx0.302\)
& \(\mathrm{QRead}(r_\Omega,s_{\rm now},\beta_{\rm sym}^2)\) with \(r_\Omega=\langle 3,\ell_{\rm now},[0,1],R_\Omega\rangle\)
& validity restricted to the present epoch, carried by the layer \(\ell_{\rm now}\)\\
\bottomrule
\end{longtable}
\end{center}

Rows~4 and~7 show where the QAM form has paid off in the project.  Several errors found in the Companion were comparisons of
quantities from different frames, or values read as derived when they were put in.  In the QAM form such a step is not
well-formed.


\begin{thebibliography}{9}
\bibitem{companion} M.~Brendecke, \emph{Quantum Sphaera Companion: A Structure-First Mathematical Unfolding},
Zenodo, concept DOI \href{https://doi.org/10.5281/zenodo.19210728}{10.5281/zenodo.19210728}.
\bibitem{bialynicki} I.~Bialynicki-Birula, Photon wave function, \emph{Prog.\ Opt.} \textbf{36} (1996) 245--294.
\bibitem{baez} J.~C.~Baez, The octonions, \emph{Bull.\ Amer.\ Math.\ Soc.} \textbf{39} (2002) 145--205.
\bibitem{hopf1931} H.~Hopf, \"Uber die Abbildungen der dreidimensionalen Sph\"are auf die Kugelfl\"ache, \emph{Math.\ Ann.} \textbf{104} (1931) 637--665.
\bibitem{hopf1935} H.~Hopf, \"Uber die Abbildungen von Sph\"aren auf Sph\"aren niedrigerer Dimension, \emph{Fund.\ Math.} \textbf{25} (1935) 427--440.
% TODO: Dixon; Furey; Gresnigt; Koide; Brannen; Ashtekar--Baez--Corichi--Krasnov; Domagala--Lewandowski; Meissner; Cartan--Kaehler.
\end{thebibliography}
\end{document}
